% Generated draft: requires mathematical and source review.
\documentclass[11pt]{article}
\usepackage[margin=25mm]{geometry}
\usepackage{fontspec}
\setmainfont{Latin Modern Roman}
\usepackage{amsmath,amssymb,mathtools,unicode-math}
\setmathfont{Latin Modern Math}
\usepackage{enumitem,needspace,xurl,hyperref}
\setlength{\parindent}{0pt}
\setlength{\parskip}{5pt}
\setlength{\emergencystretch}{3em}
\newcommand{\N}{\mathbb{N}}
\newcommand{\Z}{\mathbb{Z}}
\newcommand{\Q}{\mathbb{Q}}
\newcommand{\R}{\mathbb{R}}
\newcommand{\C}{\mathbb{C}}
\newcommand{\F}{\mathbb{F}}
\newcommand{\A}{\mathbb{A}}
\newcommand{\PP}{\mathbb{P}}
\newcommand{\E}{\mathbb{E}}
\newcommand{\eps}{\varepsilon}
\newcommand{\dd}{\,\mathrm{d}}
\begin{document}
\section*{UnsolvedMath ID: 20002044}
\textbf{Source problem code:} catalogue entry
\subsection{Definitions and mathematical statement}
Problem \(23\) asks whether a preceding detailed calculation was necessary and what it implies for an operator \(S\) from one-forms to four-forms. The excerpt supplies neither the calculation nor a definition of \(S\). The aid reconstructs its setting as an oriented conformal four-manifold, with the critical Paneitz operator in density/top-form form
\[
 P_4:\Omega^0\to\Omega^4,\qquad P_4=S_gd,
\]
and \(S_g\) an Eastwood--Slov\'ak extension of the Branson--Gover gauge companion, up to normalization and density identifications. Critical \(P_4\) is conformally invariant in this form and annihilates locally constant functions.

The precise local factorization consequence is: on the sheaf of closed one-forms, \(S_g\) is determined by \(P_4\) alone. For \(\alpha\) closed, choose local potentials \(\alpha=\dd f\), and set
\[
 S_g\alpha=P_4f.
\]
Changing \(f\) by a locally constant function changes nothing, so the local outputs patch. Critical conformal covariance then makes this restriction conformally invariant without computing an arbitrary extension on all one-forms. This uses local potentials and does not require a global potential for every closed form.

If \(S_g\) is one local differential extension, all extensions with the same factorization form the affine class
\[
 S_g+A_g,\qquad A_gd=0.
\]
They agree on closed one-form germs but need not agree on other inputs. Thus \(P_4\) does not canonically determine the full extension.

For the standard extension, the aid gives the exact defect, with \(\widehat g=e^{2\omega}g\),
\[
 S_{\widehat g}\alpha-S_g\alpha
 =-4\langle L_1\alpha,\dd\omega\rangle_g\,\dd V_g,
 \qquad L_1\alpha=\tfrac12\delta_gd\alpha,
\]
under its top-form and operator normalization. Consequently \(S\) is invariant on \(\ker L_1\), and the coupled system \((L_1,S)\) is the natural elliptically coercive gauge object in the Riemannian setting. On a closed Riemannian manifold, \(\ker(\delta d)\) consists of closed one-forms; locally the kernel condition need not be replaced by global exactness. The defect formula for a chosen extension is stronger information than the factorization alone and is not asserted for every \(S+A\).

The aid says the long component expansion is unnecessary for its Weyl-gauge invariance argument and instead verifies an already manifest curvature representative. Because the preceding gauge construction is missing, that claim is retained as methodological context rather than a complete theorem about an undefined Weyl-gauge object. The local factorization and explicit standard-extension defect are the source-specific partial statements available here.
\subsection{Short English statement}
The reconstructed Paneitz factorization determines \(S\) canonically on closed one-forms through local potentials, but leaves an affine family of extensions elsewhere. For the standard gauge companion, its conformal defect is controlled by the Maxwell operator, yielding invariance on that kernel. The missing preceding calculation prevents a broader standalone gauge statement.
\subsection{Sources}
\begin{itemize}
\item Local UnsolvedMath registry, permanent ID \(20002044\), source code \texttt{catalogue entry}; canonical collection \texttt{aim-geometry-notes.json}, notes[381].
\item AIM workshop ``Conformal structure in geometry, analysis, and physics,'' Problem \(23\): \url{https://aimath.org/WWN/confstruct/confstruct.pdf}. Exact question asks whether the detailed calculation was necessary and about \(S:\text{one-forms}\to\text{four-forms}\); the preceding definitions are not reproduced. No PDF was opened in this session.
\item The aid names Eastwood--Slov\'ak, Branson--Gover, Paneitz and the Maxwell gauge-companion interpretation, without full authoritative references. Original remarks and literature fields are empty; no proposer or year is supplied.
\item Supplied research attempt \(1\), \texttt{partial\_result}, run status \texttt{partially\_solved}, rigor label ``proved,'' novelty confidence ``low.'' Registry classification \texttt{PARTIAL-PROGRESS}, updated 14 August 2026.
\end{itemize}
\end{document}
